% ==================================================================
% id:       lim-008
% topic:    גבולות ורציפות של פונקציות
% subtopic: גבולות מופלאים
% solution: full
% ==================================================================
\textbf{
היעזרו בגבולות המופלאים שנלמדו בכיתה,
ובכל אחד מהסעיפים הבאים,
חשבו את הגבול: }
\begin{enumerate}\bMs{2}
\item $\ds{ \lim_{x \to 0} \frac{\tan(7x)}{\sin(5x)}}$

\ifsol
\begin{solbox}
נציג
\begin{align*}
    \frac{\tan(7x)}{\sin(5x)}&=\ \frac{\sin(7x)}{\cos(7x)\cdot\sin(5x)}\\
    &=\ \frac{\sin(7x)}{\color{red}7x}\cdot\frac{\color{red}7}{{\color{red}5}\cos(7x)}\cdot\frac{\color{red}5x}{\sin(5x)},
\end{align*}
לכן
\begin{align*}
    .\lim_{x\to0}\frac{\tan7x}{\sin5x}
    &=\ \lim_{x\to0}\frac{\sin(7x)}{7x}\cdot\frac{7}{5\cos(7x)}\cdot\frac{5x}{\sin(5x)}
    \overset{(\ast)}{=}\ 1\cdot\frac{7}{5}\cdot1
    =\frac{7}{5}
\end{align*}
$(\ast)$
יש כאן שימוש פעמיים בגבול מופלא.
הערה: בחישוב הגבול השתמשנו באריתמטיקה של גבולות: הגבול של סכום, של מכפלה ושל מנה של פונקציות מתכנסות שווה לסכום, למכפלה ולמנה של הגבולות, כל עוד הגבול של המכנה שונה מ-$0$.
\end{solbox}
\fi

\item $\ds{ \lim_{x \to 0} \frac{1-\cos(16x)}{\tan(5x)\sin(3x)}}$

\ifsol
\hint{השתמשו בגבול $\lim_{t\to0}\frac{1-\cos t}{t^2}$ (או בזהות $1-\cos t=2\sin^2\frac t2$), והכפילו וחלקו כך שיופיעו הגבולות המופלאים.}
\begin{solbox}
נזכר בגבול מופלא
$$.\lim_{x\to0}\frac{1-\cos x}{x^2}=\frac{1}{2}$$
נציג
\begin{align*}
    \frac{1-\cos(16x)}{\tan(5x)\sin(3x)}&=\ \frac{1-\cos(16x)}{\color{red}(16x)^2}\cdot\frac{1}{\cos(5x)}\cdot\frac{\color{red}5x}{\sin(5x)}\cdot\frac{\color{red}3x}{\sin(3x)}{\color{red}\cdot\frac{16^2}{5\cdot3}}
\end{align*}
מכאן
\begin{align*}
    \lim_{x\to0}\frac{1-\cos16x}{\tan(5x)\sin(3x)}
    &=\ \frac{16^2}{5\cdot3}\cdot\lim_{x\to0} \frac{1-\cos(16x)}{(16x)^2}\cdot\frac{1}{\cos(5x)}\cdot\frac{5x}{\sin(5x)}\cdot\frac{3x}{\sin(3x)}\\
    &=\ \frac{256}{15}\cdot\frac{1}{2}\cdot1\cdot1\cdot1
    = \frac{128}{15}.
\end{align*}

נציג עוד פתרון. נזכר בזהות
$$.1-\cos x=2\sin^2(\tfrac{x}{2})$$
נציג
\begin{align*}
    \frac{1-\cos16x}{\tan(5x)\sin(3x)}&=\frac{2\sin^2(8x)}{\tan(5x)\sin(3x)}
    = \frac{2\sin(8x)\sin(8x)\cos(5x)}{\sin(5x)\sin(3x)}\\
    &=\ 2\cdot\frac{\sin(8x)}{\color{red}8x}\cdot\frac{\sin(8x)}{\color{red}8x}\cdot\frac{\color{red}5x}{\sin(5x)}\cdot\frac{\color{red}3x}{\sin(3x)}\cdot\frac{{\color{red}8\cdot8}\cdot\cos(5x)}{\color{red}5\cdot3}\\
    &=\ \frac{128}{15}\cdot\cos(5x)\cdot\frac{\sin(8x)}{8x}\cdot\frac{\sin(8x)}{8x}\cdot\frac{5x}{\sin(5x)}\cdot\frac{3x}{\sin(3x)},
\end{align*}
מכאן
\begin{align*}
    \lim_{x\to0}\frac{1-\cos(16x)}{\tan(5x)\sin(3x)}
    &=\ \lim_{x\to0} \frac{128}{15}\cdot\cos(5x)\cdot\frac{\sin(8x)}{8x}\cdot\frac{\sin(8x)}{8x}\cdot\frac{5x}{\sin(5x)}\cdot\frac{3x}{\sin(3x)}\\
    &= \frac{128}{15}\cdot1\cdot1\cdot1\cdot1\cdot1
    =\frac{128}{15}.
\end{align*}
הערה: בחישוב הגבול השתמשנו באריתמטיקה של גבולות: הגבול של סכום, של מכפלה ושל מנה של פונקציות מתכנסות שווה לסכום, למכפלה ולמנה של הגבולות, כל עוד הגבול של המכנה שונה מ-$0$.
\end{solbox}
\fi

\item (*) $\ds{ \lim_{x \to {\pi}} \frac{\sin^2 x}{1+ \cos^3 x}}$

\ifsol
\hint{הציבו $t=x-\pi$ כדי לעבור לגבול ב-$0$, ופרקו את $1-\cos^3t$ כהפרש חזקות שלישיות.}
\begin{solbox}
נציג
\begin{align*}
    \frac{\sin^2 x}{1+ \cos^3 x}&=\ \frac{\sin^2(x-\pi+\pi)}{1+ \cos^3(x-\pi+\pi)}\\
    &=\ \frac{\sin^2(x-\pi)}{1-\cos^3(x-\pi)}.
\end{align*}
נסמן
$t=x-\pi$.
כאשר
$x$
שואף ל-
$\pi$,
$t$
שואף לאפס.
נמשיך להציג את הפונקציה בצורה נוחה יותר לטיפול:
\begin{align*}
    \frac{\sin^2(x-\pi)}{1-\cos^3(x-\pi)}=\frac{\sin^2t}{1-\cos^3t}
    = \frac{\sin^2t}{\color{red}t^2}\cdot\frac{\color{red}t^2}{1-\cos t}\cdot\frac{1}{\cos^2t+\cos t+1}
\end{align*}
נשים לב כי
כאשר
$x$
שואף ל-
$\pi$,
$t$
שואף לאפס,
מכאן
\begin{align*}
    \lim_{x \to {\pi}} \frac{\sin^2 x}{1+ \cos^3 x}&=\ \lim_{t\to0}\frac{\sin^2t}{\color{red}t^2}\cdot\frac{\color{red}t^2}{1-\cos t}\cdot\frac{1}{\cos^2t+\cos t+1}\\
    = 1\cdot2\cdot\frac{1}{3}
    = \frac{2}{3}.
\end{align*}
הערה: בחישוב הגבול השתמשנו באריתמטיקה של גבולות: הגבול של סכום, של מכפלה ושל מנה של פונקציות מתכנסות שווה לסכום, למכפלה ולמנה של הגבולות, כל עוד הגבול של המכנה שונה מ-$0$.
\end{solbox}
\fi

\item $\ds{ \lim_{x \to 0} \frac{\ln(1+3x)}{1 - \sqrt{1 + \tan x}}}$

\ifsol
\hint{הכפילו וחלקו כך שיופיעו הגבולות המופלאים של $\ln(1+t)$ ושל $(1+t)^{1/2}$.}
\begin{solbox}
נציג
\begin{align*}
    \frac{\ln(1+3x)}{1 - \sqrt{1 + \tan x}}&=\ \frac{\ln(1+3x)}{\color{red}3x}\cdot\frac{\color{red}\tan x}{1-(1+\tan x)^{\frac{1}{2}}}\cdot{\color{red}\frac{x}{\tan x}\cdot3}\\
    &=\ \frac{\ln(1+3x)}{3x}\cdot\frac{\tan x}{1-(1+\tan x)^{\frac{1}{2}}}\cdot{\frac{x}{\sin x}\cdot3\cdot\cos x}
\end{align*}
לכן
\begin{align*}
    \lim_{x\to0}\frac{\ln(1+3x)}{1 - \sqrt{1 + \tan x}}
    &=\ 3\cdot\frac{\ln(1+3x)}{3x}\cdot\frac{\tan x}{1-(1+\tan x)^{\frac{1}{2}}}\cdot{\frac{x}{\sin x}\cdot\cos x}\\
    &=\ -3\cdot1\cdot(-2)\cdot1\cdot1\\
    &=\ -6.
\end{align*}
כאשר נעזרנו בגבול המופלא
$$.\lim_{t\to0}
\frac{(1+t)^{1/2}-1}{t}=\frac12$$
הערה: בחישוב הגבול השתמשנו באריתמטיקה של גבולות: הגבול של סכום, של מכפלה ושל מנה של פונקציות מתכנסות שווה לסכום, למכפלה ולמנה של הגבולות, כל עוד הגבול של המכנה שונה מ-$0$.
\end{solbox}
\fi

\item $\ds{ \lim_{x \to 0} \frac{3^{2x} - 1}{x}}$

\ifsol
\begin{solbox}
מתקיים
\begin{align*}
    \lim_{x \to 0} \frac{3^{2x} - 1}{x}={\color{red}2}\lim_{x \to 0} \frac{3^{2x} - 1}{{\color{red}2\cdot}x}
    = 2\ln3
\end{align*}
כאשר נעזרנו בגבול המופלא
$$.\lim_{t\to0}\frac{a^t-1}{t}=\ln(a)$$
הערה: בחישוב הגבול השתמשנו באריתמטיקה של גבולות: הגבול של מכפלה של פונקציות מתכנסות שווה למכפלת הגבולות.
\end{solbox}
\fi

\item $\ds{ \lim_{x \to 0} \sqrt[x]{1-3x}}$

\ifsol
\hint{זהו גבול מהצורה $1^\infty$: התאימו את המעריך כדי להשתמש בגבול $\lim_{t\to0}(1+t)^{1/t}=e$.}
\begin{solbox}
נציג
\begin{align*}
    \sqrt[x]{1-3x}&=\ (1-3x)^{\tfrac{1}{x}}\\
    &=\ (1-3x)^{\tfrac{1}{{\color{red}-3}x}\cdot{\color{red}(-3)}},
\end{align*}
מכאן
\begin{align*}
    \lim_{x \to 0} \sqrt[x]{1-3x}&=\ \lim_{x \to 0}\left((1-3x)^{\tfrac{1}{{-3}x}}\right)^{-3}\\
    &=\ e^{-3}.
\end{align*}
כאשר נעזרנו בגבול המופלא
$$.\lim_{t\to0}(1+t)^{1/t}=e$$
\end{solbox}
\fi

\item $\ds{ \lim_{x \to 0} \frac{e^{ax}-e^{bx}}{x}}$

\ifsol
\hint{הוציאו את $e^{bx}$ כגורם משותף מהמונה.}
\begin{solbox}
אם
$a=b$,
אז הגבול הוא
$0$.
נניח כי
$a\neq b$,
נציג
\begin{align*}
    \frac{e^{ax}-e^{bx}}{x}&=\ e^{bx}\cdot\frac{e^{(a-b)x}-1}{x}\\
    &=\ e^{bx}\cdot\frac{e^{(a-b)x}-1}{{\color{red}(a-b)}x}\cdot{\color{red}(a-b)},
\end{align*}
מכאן
\begin{align*}
    \lim_{x\to 0}\frac{e^{ax}-e^{bx}}{x}
    &=\ (a-b)\cdot\lim_{x\to 0}e^{bx}\cdot\frac{e^{(a-b)x}-1}{{(a-b)}x}\\
    &=\ (a-b)\cdot1\cdot1\\
    &=\ a-b.
\end{align*}
נעזרנו בגבול מופלא שכבר ציטטנו בסעיף קודם (מהו?).
נשים לב כי התוצאה שקיבלנו מכילה גם את המקרה שבו
$a=b$.
הערה: בחישוב הגבול השתמשנו באריתמטיקה של גבולות: הגבול של מכפלה של פונקציות מתכנסות שווה למכפלת הגבולות.
\end{solbox}
\fi

\item $\ds{ \lim_{x \to 0} (\cos(4x))^{\frac{1}{\cos x - \cos(7x)}}}$

\ifsol
\hint{זהו גבול מהצורה $1^\infty$. כתבו את הבסיס כ-$1+$(ביטוי ששואף לאפס), וחשבו את גבול המעריך בעזרת הגבול של $\frac{1-\cos t}{t^2}$.}
\begin{solbox}
בגלל שכאשר
$x\to0$
נקבל מצב של
$"1^{1/0}"$,
ננסה להיעזר בגבול המופלא
$$\lim_{t\to0}(1+t)^{1/t}=e$$
ולכן
נציג
\begin{align*}
    (\cos(4x))^{\frac{1}{\cos x - \cos(7x)}}&=\ ({\color{red}1}+\cos(4x){\color{red}-1})^{\frac{1}{\cos x - \cos(7x)}}\\
    &=\ (1+(\cos(4x)-1))^{{\color{red}\frac{1}{\cos(4x)-1}\cdot}\frac{\color{red}\cos(4x)-1}{\cos x - \cos(7x)}}.
\end{align*}
נטפל במעריך
\begin{align*}
    \frac{\cos(4x)-1}{\cos x - \cos(7x)}&=\frac{\cos(4x)-1}{\color{red}(4x)^2}\cdot\frac{\color{red}(4x^2)}{\cos x{\color{red}-1+1}-\cos(7x)}\\
    &=\ \frac{\cos(4x)-1}{(4x)^2}\cdot\frac{1}{\frac{1}{16}\cdot\frac{\cos x-1}{x^2}+\frac{\color{red}49}{16}\frac{1-\cos(7x)}{{\color{red}7^2}\cdot x^2}}
\end{align*}
מתקיים
\begin{align*}
\lim_{x\to0}\frac{\cos(4x)-1}{\cos x - \cos(7x)}&=\ \lim_{x\to0}\frac{\cos(4x)-1}{(4x)^2}\cdot\frac{1}{\frac{1}{16}\cdot\frac{\cos x-1}{x^2}+\frac{\color{red}49}{16}\frac{1-\cos(7x)}{{\color{red}7^2}\cdot x^2}}\\
&=\ -\frac{1}{2}\cdot\frac{1}{\frac{1}{16}\cdot\left(-\frac{1}{2}\right)+\frac{49}{16}\cdot\frac{1}{2}}\\
&=\ -\frac{1}{3}
\end{align*}
לכן
בעזרת הגבול המופלא שהזכרנו בתחילת הפתרון, נקבל
    $$\lim_{x \to 0} (\cos4x)^{\frac{1}{\cos x - \cos 7x}}=e^{-\frac{1}{3}}$$
\end{solbox}
\fi

\item $\ds{ \lim_{x \to -3} (x^2 - 8)^{\frac{1}{x^2 + x - 6}}}$

\ifsol
\hint{זהו גבול מהצורה $1^\infty$. כתבו את הבסיס כ-$1+$(ביטוי ששואף לאפס), ופרקו לגורמים את מה שמתקבל במעריך.}
\begin{solbox}
מאותו השיקול של הסעיף הקודם,
נציג
\begin{align*}
    (x^2 - 8)^{\frac{1}{x^2 + x - 6}}&=\ (1+(x^2-9))^{{\color{red}\frac{1}{x^2-9}}\cdot\frac{\color{red}x^2-9}{x^2+x-6}}.
\end{align*}
נשים לב כי
$$,\frac{x^2-9}{x^2+x-6}=\frac{(x-3)(x+3)}{(x-2)(x+3)}=\frac{x-3}{x-2}$$
לכן
$$.\lim_{x\to-3}\frac{x^2-9}{x^2+x-6}=\lim_{x\to-3}\frac{x-3}{x-2}=\frac{-3-3}{-3-2}=\frac{6}{5}$$
מתקיים
\begin{align*}
    \lim_{x \to -3} (x^2 - 8)^{\frac{1}{x^2 + x - 6}}&=\ \lim_{x \to -3}(1+(x^2-9))^{\frac{1}{x^2-9}\cdot\frac{x-3}{x+2}}\\
    &=\ e^{\frac{6}{5}}.
\end{align*}
\end{solbox}
\fi

\item
$\ds{
\lim_{x\to 0}\frac{\tan(x) - \sin(x)}{\left( \sin(x) \right) ^3}
}$

\ifsol
\hint{כתבו $\tan x=\frac{\sin x}{\cos x}$ והוציאו גורם משותף במונה.}
\begin{solbox}
נציג
\begin{align*}
    \frac{\tan(x) - \sin(x)}{\left( \sin(x) \right) ^3}&=\ \frac{\frac{\sin x}{\cos x}-\sin x}{\sin^3x}\\
    &=\ \frac{\sin x(1-\cos x)}{\sin^3x\cdot\cos x}\\
    &=\ \frac{1-\cos x}{\sin^2x\cdot\cos x}\\
    &=\ \frac{1-\cos x}{\color{red}x^2}\cdot\frac{\color{red}x^2}{\sin^2x}\cdot\frac{1}{\cdot\cos x}
\end{align*}
מכאן
בעזרת אריתמטיקה של כפל גבולות ניעזר בגבולות המופלאים המוכרים ונקבל
\begin{align*}
    \lim_{x\to0}\frac{\tan(x) - \sin(x)}{\left( \sin(x) \right) ^3}&=\ \lim_{x\to0}\frac{1-\cos x}{\color{red}x^2}\cdot\frac{\color{red}x^2}{\sin^2x}\cdot\frac{1}{\cdot\cos x}\\
    &=\ \frac{1}{2}\cdot1\cdot1\\
    &=\ \frac{1}{2}.
\end{align*}
\end{solbox}
\fi

\item $\ds{ \lim_{x \to \infty} (2x+3)\lr{\ln(x + 2) - \ln(x)}}$

\ifsol
\hint{השתמשו בחוקי הלוגריתמים כדי לכתוב את הביטוי כ-$\ln$ של חזקה, ושימו לב שבתוך ה-$\ln$ יש גבול מהצורה $1^\infty$.}
\begin{solbox}
ניעזר בחוקי לוגריתמים ונקבל
\begin{align*}
    (2x+3)\lr{\ln(x + 2) - \ln(x)}=(2x+3)\ln(1+\tfrac{2}{x})
    = \ln\left((1+\tfrac{2}{x})^{2x+3}\right)
    \end{align*}
    הביטוי בתוך ה-$\ln$ דומה לגבול המופלא המוכר, ולכן נמשיך לפשט רק אותו  ה-$\ln$-
    \begin{align*}
    (1+\tfrac{2}{x})^{2x+3}
    &=\ (1+\tfrac{2}{x})^{{\color{red}\frac{x}{2}}\cdot(2x+3)\cdot{\color{red}\frac{2}{x}}}\\
    &=\ (1+\tfrac{2}{x})^{{\frac{x}{2}}\cdot{(4+\tfrac{6}{x})}}
\end{align*}
לכן
\begin{align*}
    \lim_{x\to\infty}(2x+3)\lr{\ln(x + 2) - \ln(x)}&=\ \lim_{x\to\infty}(4+\tfrac{6}{x})\ln(1+\tfrac{2}{x})^{\frac{x}{2}}\\
    &=\ 4\cdot\ln e
    = 4.
\end{align*}
הערה: בחישוב הגבול השתמשנו באריתמטיקה של גבולות: הגבול של מכפלה של פונקציות מתכנסות שווה למכפלת הגבולות.
\end{solbox}
\fi

\item
$\ds{
\lim_{x\to 1} (1 - x) \cdot \tan\lr{\frac\pi2 x}
}$

\ifsol
\hint{הציבו $t=1-x$ והשתמשו בזהות $\tan(\frac\pi2-\theta)=\cot\theta$.}
\begin{solbox}
נסמן
$t=1-x$.
לכן
$$,(1 - x) \cdot \tan\lr{\frac\pi2 x}=t\cdot\tan\left(\frac{\pi}{2}-\frac{\pi}{2}t\right)=t\cdot\cot\left(\frac{\pi}{2}t\right)=\frac{t}{\sin\left(\frac{\pi}{2}t\right)}\cdot\cos\left(\frac{\pi}{2}t\right)$$
מכיוון ש-
$x\longrightarrow1$,
נסיק כי
$t\longrightarrow0$,
לכן
\begin{align*}
    \lim_{x\to 1} (1 - x) \cdot \tan\lr{\frac\pi2 x}&=\ \lim_{t\to 0}\frac{t}{\sin\left(\frac{\pi}{2}t\right)}\cdot\cos\left(\frac{\pi}{2}t\right)\\
    &=\ {\color{red}\frac{2}{\pi}\cdot}\lim_{t\to 0}\frac{{\color{red}\frac{\pi}{2}}t}{\sin\left(\frac{\pi}{2}t\right)}\cdot\cos\left(\frac{\pi}{2}t\right)\\
    &=\ \frac{2}{\pi}\cdot1\cdot1
    = \frac{2}{\pi}.
\end{align*}
הערה: בחישוב הגבול השתמשנו באריתמטיקה של גבולות: הגבול של מכפלה של פונקציות מתכנסות שווה למכפלת הגבולות.
\end{solbox}
\fi

\item
$\ds{
\lim_{x\to 0}\frac{2 \cdot \arcsin(x)}{3x}
}$

\ifsol
\hint{הציבו $y=\arcsin x$.}
\begin{solbox}
נסמן
$y=\arcsin x$,
מכאן
$x=\sin y$.
נשים לב כי כאשר
$x\to 0$,
גם
$y\to 0$.
מתקיים
$$,\frac{2 \cdot \arcsin(x)}{3x}=\frac{2y}{3\sin y}$$
מכאן
\begin{align*}
    \lim_{x\to0}\frac{2 \cdot \arcsin(x)}{3x}= \lim_{y\to0}\frac{2y}{3\sin y}
    = \frac{2}{3}.
\end{align*}
הערה: בחישוב הגבול השתמשנו באריתמטיקה של גבולות: הגבול של סכום, של מכפלה ושל מנה של פונקציות מתכנסות שווה לסכום, למכפלה ולמנה של הגבולות, כל עוד הגבול של המכנה שונה מ-$0$.
\end{solbox}
\fi

\item $\ds{\lim_{x\to0}\frac{\sqrt[5]{1+2x}-\sqrt[5]{1+x}}{x}}$

\ifsol
\hint{הוסיפו והחסירו $1$ במונה, והשתמשו בגבול $\lim_{t\to0}\frac{(1+t)^a-1}{t}=a$.}
\begin{solbox}
נציג
\begin{align*}
    \frac{\sqrt[5]{1+2x}-\sqrt[5]{1+x}}{x}
    &=\ \frac{(1+2x)^{\frac15}-1}{x}-\frac{(1+x)^{\frac15}-1}{x}\\
    &=\ {\color{red}2}\cdot\frac{(1+2x)^{\frac15}-1}{{\color{red}2}x}-\frac{(1+x)^{\frac15}-1}{x}.
\end{align*}
לכן לפי הגבול המופלא
$$\lim_{t\to0}\frac{(1+t)^a-1}{t}=a$$
נקבל
\begin{align*}
    \lim_{x\to0}\frac{\sqrt[5]{1+2x}-\sqrt[5]{1+x}}{x}
    &=\ 2\cdot\frac15-\frac15\\
    &=\ \frac15.
\end{align*}
הערה: בחישוב הגבול השתמשנו באריתמטיקה של גבולות: הגבול של סכום, של מכפלה ושל מנה של פונקציות מתכנסות שווה לסכום, למכפלה ולמנה של הגבולות, כל עוד הגבול של המכנה שונה מ-$0$.
\end{solbox}
\fi
\eMs
\end{enumerate}
